\akhead{Topic 4 — Systems of Two Linear Equations in Two Variables}

\ans{1}{1}{$x=7$}{Add the two equations: $2x=14$, so $x=7$ (and $y=3$).}

\ans{1}{2}{C}{Substitute $y=3x$ into $x+y=16$: $4x=16$, so $x=4$ and
$y=3(4)=12$. Choice A is $x$, not $y$.}

\ans{1}{3}{B}{Subtract the first equation from the second: $x=5$. Then
$5+y=9$, so $y=4$. Choice C is $x$; choice A comes from subtracting in the
wrong order.}

\ans{1}{4}{$y=5$}{Substitute $y=x+1$: $2x+3(x+1)=23$, so $5x=20$ and $x=4$.
Then $y=4+1=5$.}

\ans{1}{5}{B}{Add the equations: $4x=12$, so $x=3$. (Then $3y=6$ and $y=2$,
which is choice A.)}

\ans{1}{6}{B}{With $s$ small and $l$ large boxes, $s+l=9$ and $2s+5l=33$.
Substituting $s=9-l$ gives $18+3l=33$, so $l=5$ (and $s=4$, choice A).}

\ans{1}{7}{B}{Substitute: $5x-(x+2)=22$, so $4x=24$ and $x=6$. Choice A
comes from dropping the sign on the $2$; choice D stops at $4x=24$; choice C
is $y$.}

\ans{1}{8}{C}{Add the equations: $3x=9$, so $x=3$, and then $y=2$. Both
equations check: $2(3)+2=8$ and $3-2=1$. Choice B reverses the coordinates.}

\ans{1}{9}{$y=2$}{Add the equations: $8x=40$, so $x=5$. Then $15+2y=19$
gives $y=2$.}

\ans{1}{10}{B}{From the second equation $y=2x-5$. Then $4x+3(2x-5)=25$, so
$10x=40$ and $x=4$. Choice A is $y$; choice D is $x+y$.}

\ans{1}{11}{C}{Add the two equations: $7x+7y=49$, so $x+y=7$ with no need to
find $x$ and $y$. Choice A is $x-y$, choice B is $x$, choice D forgets to
divide by $7$.}

\ans{1}{12}{$65$}{With $c$ and $t$ the two costs, $3c+2t=430$ and $c+t=165$.
Substituting $c=165-t$ gives $495-t=430$, so $t=65$ (and $c=100$).}

\ans{1}{13}{C}{With $a+s=240$ and $12a+7s=2180$, substitute $a=240-s$:
$2880-5s=2180$, so $5s=700$ and $s=140$. Choice A is the number of adult
tickets.}

\ans{1}{14}{A}{The equation $x+y=45$ adds the hours on the two tasks, so
$45$ is the total number of hours worked. The dollar total $440$ appears in
the second equation.}

\ans{1}{15}{B}{Multiply the first equation by $2$: $x+2y=14$. With $x=y+2$,
$3y+2=14$, so $y=4$ and $x=6$. Choice C comes from writing $x+y=14$ instead
of $x+2y=14$.}

\ans{1}{16}{B}{Multiply the second equation by $2$: $10x-4y=28$. Adding
gives $13x=26$, so $x=2$ and $3(2)+4y=-2$ gives $y=-2$. Choice C is $x$.}

\ans{1}{17}{$5$}{Set the costs equal: $30+4h=10+8h$, so $20=4h$ and $h=5$
hours.}

\ans{1}{18}{$k=4$}{The pair $(3,2)$ satisfies the first equation. Put it in
the second: $4(3)+k(2)=20$, so $2k=8$ and $k=4$.}

\ans{1}{19}{B}{Divide the first equation by $3$: $x+2y=5$. The lines are
parallel when the $y$-coefficients match, so $k=2$; the constants $5$ and
$9$ differ, so there is no solution. Choice D copies $6$ without scaling.}

\ans{1}{20}{C}{Multiply the first equation by $-\tfrac{3}{2}$:
$-6x+3y=-15$, which is the second equation exactly when $a=3$. Choice A
drops the sign of the multiplier.}

\ans{1}{21}{$3$}{Subtract the first equation from the second:
$-2x+2y=6$, so $y-x=3$ without solving for either variable.}

\ans{1}{22}{C}{Two lines in slope-intercept form are parallel when the
slopes match and the intercepts do not. Here $c=4$ gives slope $4$ with
intercepts $7$ and $-3$, so there is no solution.}

\ans{2}{1}{$x=4$}{Substitute $y=13-3x$ into $2x-3y=5$: $2x-39+9x=5$, so
$11x=44$ and $x=4$ (with $y=1$).}

\ans{2}{2}{B}{Add the equations: $10x=30$, so $x=3$. Then $18+5y=28$ gives
$y=2$. Choice C is $x$.}

\ans{2}{3}{C}{With $l+w=23$ from the perimeter and $l=w+5$, $2w+5=23$, so
$w=9$ and $l=14$. Choice A is the width and choice D is $l+w$.}

\ans{2}{4}{$2$}{Subtract the second equation from the first:
$4x+4y=8$, so $x+y=2$ directly.}

\ans{2}{5}{A}{Add the equations: $13x+13y=65$, so $x+y=5$. Choice B is
$x-y$ and choice D forgets to divide by $13$.}

\ans{2}{6}{B}{With $n+d=30$ and $5n+10d=220$, divide the second by $5$:
$n+2d=44$. Subtracting gives $d=14$ (and $n=16$, choice C).}

\ans{2}{7}{B}{The pair $(3,2)$ satisfies $2x-y=4$. Put it in the first
equation: $3a+6=18$, so $3a=12$ and $a=4$. Choice D forgets to divide by
$3$.}

\ans{2}{8}{$16$}{Subtract the second equation from the first: $12y=24$, so
$y=2$; then $4x=14+6=20$ and $x=5$. Hence $2x+3y=10+6=16$.}

\ans{2}{9}{B}{Add $2y$ to both sides of $x-2y=7$ to get $x=2y+7$. Choice A
subtracts instead of adds; choice C moves $x$ rather than $2y$.}

\ans{2}{10}{B}{The equation $3n+2f=96$ multiplies each count by its price,
so $96$ is the total cost in dollars. The count total $40$ is the first
equation.}

\ans{2}{11}{$35$}{With $a+b=60$ and $12a+18b=870$, substitute $b=60-a$:
$1080-6a=870$, so $6a=210$ and $a=35$ minutes.}

\ans{2}{12}{B}{Double both equations: $x+2y=18$ and $2x-y=16$. Then
$y=2x-16$, so $x+4x-32=18$, giving $x=10$ and $y=4$. Choice D is $x$.}

\ans{2}{13}{D}{Multiplying the second equation by $-2$ gives $4x-6y=10$, the
first equation. The two equations describe the same line, so every point on
it is a solution.}

\ans{2}{14}{$5$}{Multiply the first equation by $2$: $10x+4y=8$. Adding
gives $13x=26$, so $x=2$ and $10+2y=4$ gives $y=-3$. Then $x-y=5$.}

\ans{2}{15}{C}{With $y=2$, the second equation gives $x-6=-4$, so $x=2$.
Then $2k+8=20$, so $k=6$. Choice D drops the $4y$ term.}

\ans{2}{16}{$11$}{Both equations reduce to $2x+3y=11$: divide the first by
$3$ and the second by $2$. That is why the system has infinitely many
solutions, and $2x+3y=11$ for every one of them.}

\ans{2}{17}{A}{Parallel lines need $\dfrac{k}{2}=\dfrac{-6}{3}=-2$, so
$k=-4$. Then the first equation is $2x+3y=-9$, which contradicts
$2x+3y=7$, so there is no solution. Choice D drops the sign.}

\ans{2}{18}{C}{The second equation must be $3$ times the first, since
$-15=3(-5)$. So $a=3(3)=9$ and $b=3(8)=24$, giving $a+b=33$. Choice A
leaves the constant unscaled.}

\ans{2}{19}{$c=7$}{No solution requires equal ratios of the
$x$- and $y$-coefficients: $\dfrac{4}{12}=\dfrac{c}{21}$, so $c=7$. The
constants then disagree ($9$ versus $\tfrac{5}{3}$), confirming no
solution.}

\ans{2}{20}{C}{The second equation is $3$ times the first, so the system is
one line written twice; every pair satisfying $3x+2y=7$ satisfies both.}

\ans{2}{21}{A}{With $b+f=300$ and $15b+25f=6100$, substitute $b=300-f$:
$4500+10f=6100$, so $f=160$ and $b=140$. The difference is $160-140=20$.}

\ans{2}{22}{C}{Parallel lines need $\dfrac{k}{4}=\dfrac{k+1}{6}$, so
$6k=4k+4$ and $k=2$. Then $2x+3y=12$ while the second equation gives
$2x+3y=4.5$, so there is no solution.}
